% QUESTAO
Após obter a altura do retângulo através do \textbf{teorema de Pitágoras}, o cálculo de sua área, em metros quadrados, fica:

\parbox{0.4\linewidth}{

% ALTERNATIVAS
$120$\,m$^2$
$100$\,m$^2$
$60$\,m$^2$
$8$\,m$^2$
$150$\,m$^2$

}\parbox{0.59\linewidth}{
\begin{center}
\vspace{-6mm}
\begin{tikzpicture}[ >=latex, scale=0.85 ]

\retangulo{$15$}{15}{}{8};
\draw[dashed,nodelabel={8pt}{$17$}] (0,0) -- (15,8);

\end{tikzpicture}
\vspace{-9mm}
\end{center}
}


